\documentclass[11pt]{article}
\usepackage{amsmath,amssymb,amsthm}
\usepackage[margin=1in]{geometry}
\newtheorem{proposition}{Proposition}
\newcommand{\gen}[1]{\langle #1\rangle}
\title{A proposed rank-four application for position 2739}
\author{Working mathematical application; not a Root adoption}
\date{}
\begin{document}
\maketitle

\section*{Subject and status}

This application concerns only original job $[20,4]$, orbit position $2739$,
with retained row identity \texttt{job0020:position2739:J08349}.
Its canonical marking hash is
\begin{center}\small\texttt{80c22b9b749cda7e4f79b6956d827f753243aeffeff90f2e2492d0ffcb6ff843}\end{center}
and its model hash is
\begin{center}\small\texttt{252598fb2be23480f4e91a83a6d8e7818f203deb62892b7e42157d7f180538dc}.\end{center}
The corresponding original orbit row has hash
\begin{center}\small\texttt{0c9b33e90274a8f9d1c83830201346a48249eab78b47de2d31cf7d22afd2b4d2}.\end{center}
The associated representative index is $8349$.

Write the four ordered marks as $(a,b,c,d)$.  The historical generation
Source writes this same ordered tuple as $(a,b,c,e)$: its final symbol $e$
is the present $d$, not an additional generator.  The prefix export
declares selected labels $[1,4,2]$ and end order $[1,4,2,3]$.  The checker
reads its three exported matrices as $(a,e,b)$ and constructs the final
ordered quartet as $(a,b,c,e)$; this explicit reordering is retained.
All statements below concern the same
faithful right-row model over $\mathbb F_{11}$.

The machine-readable companion is a proposal with Root authority false,
adopted false, issued false, and zero canonical or new case credit.  The
abstract theorem has a distinct proof review.  This case application
still requires its own distinct review and Root's logical issue before
any disposition is changed.  No Original matrix, vector, form or group
word is decoded or evaluated in preparing this application.

\section*{Imported ambient and generation premises}

The ambient group is the finite split Chevalley group $G=G_2(11)$,
without a covering, outer automorphism or semilinear extension.  Its
nonabelian simplicity is an established group-theoretic premise already
used in the retained generation theory.  The named group and the
identification of its acquired seven-dimensional natural Cayley module
are imported premises, not conclusions inferred from a matrix hash.

The positive old frame import, with hash beginning \texttt{40838480},
owns the same four Original marks in the named $G_2(11)$ and the same
faithful frame.  Its exact subject is carried into the new transport
import, hash beginning \texttt{acd4270f}, through the duties concerning
the same marks/model/frame and the accepted old Source/actual/frame and
Original $G_2$ premises.  The literal transport metadata pins the
retained subset, the complete original domain and the same orbit row.
These ownership records are existing positive imports; this proposal
does not create replacements for them.

Generation is imported from the adopted record with full SHA256
\begin{center}\small\texttt{9ac047f77751fdfa8153caf1820bc5249654350ed1658cee4cae05310dae5472}.\end{center}
Its schema is \texttt{ROOT\_ACCEPTED\_Q11\_QUARTET37\_GENERATION\_FACTS\_V1},
its job is $[20,4]$, and position $2739$ belongs to its verified generation
list.  Its pinned direct result is
\texttt{Q11\_QUARTET37\_DIRECT\_GENERATION\_RESULT\_V1}, with a positive
record for exactly that position.  The original proposal record has
representative $8349$ and the whole orbit-row hash displayed above.
Its original-domain descriptor bank equals the bank in the current
literal transport metadata.  These exact joins give
\begin{equation}\label{eq:generation}
                         G=\gen{a,b,c,d}.
\end{equation}

The accepted generation argument combines a full cyclic span with an
actual word of prime order $37$.  Its use of the natural Cayley module,
the rank-one operator argument, and the complete maximal-overgroup
classification remain attributed to that retained argument and theory.
The complete classification is borrowed from the sources identified
there: Kleidman's classification as checked in Nguyen's displayed list
and Wilson's book, with the stated exceptional restrictions.  None of
that classification, the prime witness or the matrix identities is
reproved or recomputed here.  Generation alone was explicitly not an
intersection-property certificate in the accepted record.

\section*{Where the involution and string premises were checked}

These premises do not rest on a new declaration in this proposal.
The accepted generation direct checker is the retained Source
\texttt{verify.py}, $33921$ bytes, with full SHA256
\begin{center}\small\texttt{9f0fa443dd62f92b432167bd9904c24a804bf3d3e50f476bc83ba448cbe2ff9f}.\end{center}
Its \texttt{context} function pins the genuine original producer and
Saved terminal records, their complete nonobservational equality, the
original labels and actual frame, the named-model premises and the
faithful-frame join.  Its \texttt{actual\_marks} function reconstructs
the fourth entry from the exact original orbit representative, forms
the ordered tuple \texttt{(a,b,c,e)}, and requires
\[
 \text{four distinct marks},\qquad
 g\ne I_7\ \text{and}\ g^2=I_7\ \text{for each mark},
\]
\[
 g_i g_j=g_j g_i\qquad(0\le i<j\le3,\ j\ge i+2).
\]
Its positive-record branch calls \texttt{actual\_marks} before accepting
the positive certificate for the position, then evaluates the supplied
word.  The verified direct-result list equals the adopted list and
contains $2739$.  Therefore the previously accepted execution supplies
\begin{equation}\label{eq:string}
 a,b,c,d\ne1,\quad a^2=b^2=c^2=d^2=1,\quad
                    [a,c]=[a,d]=[b,d]=1.
\end{equation}
The Source is inspected as text and pinned in full; its numerical
functions are not imported or executed by this application.  The use
of the accepted execution, its exact positive record and its Source
contract is an import of established facts, not an independent numerical
recheck or an inference of mathematical truth from hashes alone.

\section*{The new end-intersection certificate}

The successful new execution has result SHA256
\begin{center}\small\texttt{8caa24296cdb446b9f5127408384a0a0e190ba67b4e8d612b100f465b6064000}\end{center}
and complete proof SHA256
\begin{center}\small\texttt{7411fc3921aab4640e2b928c5fa868167dcfd43ae236666326f2f7a08b304f28}.\end{center}
The distinct actual audit, beginning \texttt{efdc37b6}, supports the
same-case fixed-line and third-form end-intersection certificate.
It closes the new Source, exact Root import and request, genuine API
frames, guard and complete IO sequence.  Its result states fixed-line
dimension one, new basis rank six, and
\begin{equation}\label{eq:end}
                \gen{a,b,c}\cap\gen{b,c,d}=\gen{b,c}.
\end{equation}

This is an imported certificate under the retained theory and historical
premises: the common commutant dimension, complete middle-dihedral
journal, faithful frame and tensor facts.  Those premises and their
attribution remain in force.  The failed transport v1 stopped while
reading floating timing metadata before loading the Original data;
the immutable parent controls review and its representation-only
successor remain retained.  No successful scientific calculation is
repeated to obtain this logical application.

\section*{The logical deduction}

The independently reviewed note \emph{Normal subgroups forced by an end
intersection}, theorem Source hash beginning \texttt{e97384b1}, proves:
if a finite group with no nontrivial solvable normal subgroup is
generated by four nonidentity string involutions, then both end
triples have the intersection property.  Its final corollary says that
the full rank-four intersection property is then equivalent to
\eqref{eq:end}.  The distinct review of that note has full SHA256
\begin{center}\small\texttt{00e4a02dfc2d41a1051e2f9a45dc8304c28e6ff8afb294c9330787d7cf9f228b}.\end{center}

\begin{proposition}
Under the imported accepted premises and exact subject joins described
above, the ordered marking at position $2739$ is a rank-four string
C-group representation of $G_2(11)$.
\end{proposition}
\begin{proof}
A nonabelian simple group has no nontrivial solvable normal subgroup.
The ambient premise and \eqref{eq:generation}--\eqref{eq:string} therefore
meet every hypothesis of the end-subgroup theorem.  Hence
$(\gen{a,b,c};a,b,c)$ and $(\gen{b,c,d};b,c,d)$ both have the full
intersection property.  Using \eqref{eq:end}, the rank-four gluing
corollary gives
\[
 \gen{r_i:i\in A}\cap\gen{r_i:i\in B}
              =\gen{r_i:i\in A\cap B}
              \qquad(A,B\subseteq\{0,1,2,3\}),
\]
where $(r_0,r_1,r_2,r_3)=(a,b,c,d)$.  The family generates $G$ by
\eqref{eq:generation}, and is a string family by \eqref{eq:string},
which is exactly the asserted representation.
\end{proof}

\section*{What has and has not been replaced}

The new contribution to this logical application is the ambient normal
subgroup obstruction and its application to the two ends.  Its dihedral
uniqueness and Klein-four factor argument extend the attributed internal
rank-three proof; the rank-four gluing criterion is elementary background
proved in the new note.  No historical originality claim is made.

This removes the need to import a separate end-rank-three intersection
assertion from the unavailable old local-review body with hash beginning
\texttt{33f56de6}.  It does not reconstruct that body, and does not
replace or weaken the complete $D_{22}$ journal, named-group membership,
generation, frame or tensor premises already imported by the accepted
end-intersection certificate.  In particular, the new abstract theorem
proves no dihedral order and no missing arithmetic identity.

The proposition is a proposed logical consequence of accepted premises.
It is not a canonical disposition, a newly issued Root assertion, a
whole-manuscript recovery, an open-conjecture resolution, a Lean proof,
or a compression measurement.  A separate case review and positive Root
logical adoption remain necessary before any such case disposition is
changed.  The machine-readable companion records all exact descriptors,
the Source locations supplying each premise and its attribution.
\end{document}
